\chapter{Hello World}

\section{To create a bash script}

First, we'll have to create a file with a .sh extention. Then we can edit it in with whatever command we want.\\

\lstinputlisting[
    language=bash, 
    caption={1-bash.sh},
    captionpos=t,
    label={lst:1-bash}, 
    mathescape=true, 
    breaklines=true]
{../1-bash/1-bash.sh}

The way bash scripts are execute is top to bottom.

\section{how to run it?}

First we'll have to change the script properties from the terminal. \\

\begin{lstlisting}[
    language=bash, 
    caption={Terminal},
    captionpos=t, 
    label={lst:1-terminal}, 
    mathescape=true,
    breaklines=true]
sudo chmod +x 1-bash.sh
\end{lstlisting}

This makes the script executable. Then we just have to execute it from the terminal:\\

\begin{lstlisting}[
    language=bash, 
    caption={Terminal},
    captionpos=t,
    label={lst:1.1-terminal}, 
    mathescape=true, 
    breaklines=true]
./1-bash.sh
\end{lstlisting}
\newpage

\section{Good practice for new scripts}

What we are using here is a "shebang". This tells the interpreter whether they should be the ones 
executing this script. Otherwise, they should call the specified interpreter to execute the script.

\lstinputlisting[
    language=bash, 
    caption={1-bash.sh},
    captionpos=t,
    label={lst:1.1-bash}, 
    mathescape=true, 
    breaklines=true]
{../1-bash/1.1-bash.sh}

\section{Comments}

You can add comments to you script using `\#' followed by text.\\

\section{Hello World!}

So to write anything out, you use the echo command. Here we'll do what we always do at the beginning
and write our hello world.\\

\lstinputlisting[
    language=bash, 
    caption={1-bash.sh},
    captionpos=t,
    label={lst:1.2-bash}, 
    mathescape=true, 
    breaklines=true]
{../1-bash/1.2-bash.sh}